%% kapa: exam=1 round=34 year=2023 period=1 form=A subject=civil_law subject_ko=민법 seq=3 number=3
%% source: Q-Net 감정평가사 기출문제 (cst00309) artlSeq=5211789 file=raw/1008/5211789/2023년도 제34회 감정평가사 1차시험 1교시.pdf
%% conversion: kordoc 4.24.51 -> tools/gen1.py (auto); review=proofread-sonnet; flags=none
\begin{question}{3}
부재자의 재산관리에 관한 설명으로 옳지 않은 것은? (다툼이 있으면 판례에 따름)
\begin{choices}
  \item 부재자로부터 재산처분권한을 수여받은 재산관리인은 그 재산을 처분함에 있어 법원의 허가를 받을 필요가 없다.
  \item 부재자가 재산관리인을 정하지 않은 경우, 부재자의 채권자는 재산관리에 필요한 처분을 명할 것을 법원에 청구할 수 있다.
  \item 법원이 선임한 재산관리인은 법원의 허가 없이 부재자의 재산에 대한 차임을 청구할 수 있다.
  \item 재산관리인의 처분행위에 대한 법원의 허가는 이미 행한 처분행위를 추인하는 방법으로 할 수 있다.
  \item 부재자가 사망한 사실이 확인되면 부재자 재산관리인 선임결정이 취소되지 않더라도 관리인의 권한은 당연히 소멸한다.
\end{choices}
\end{question}
\answer{5}
